\documentclass{article}
\usepackage{mathtools}
\usepackage{unicode-math}
\setmainfont{Times New Roman}
\setsansfont{Arial}
\setmonofont{Consolas}
\setmathfont{XITSMath-Regular.otf}[BoldFont=XITSMath-Bold.otf]
\usepackage{xeCJK}
\setCJKmainfont{SimSun}
\setCJKsansfont{SimHei}
\setCJKmonofont{SimSun}
\usepackage[version=4]{mhchem}
\usepackage{physics}
\usepackage{xcolor}
\usepackage{cancel}
\usepackage[active,tightpage]{preview}
\PreviewBorder=1pt

\begin{document}

简化字指在汉字简化过程中已经被简化了的汉字，与繁体字（正体字）相对。在不同的语言背景之下，简化字经常被俗称为简体字、简笔字、俗体字，或被蔑称为残体字。由于仍有大多数汉字并未被简化，这类汉字在中国大陆的规范名称为传承字。[1]

简体中文（网页语言代码：zh-Hans），与繁体中文相对，是使用未被简化的传承字以及简化字作为字汇的一种书面表达方式，常在中国大陆、新加坡、马来西亚、缅甸佤邦等地，以及联合国、世界银行等国际组织使用。[2]

这些简化字是根据《汉字简化方案》中的规则把传统中文中的一些汉字进行简化、合并、改造、新创而来，以及搜集在中国历史文献中已有的异体字、俗字、潦草字、书法连写、书法简写、古字而来的。

1977年颁布的二简字造成了一些混乱，更不受人们欢迎，在1986年撤回。[3]

2009年5月，中华人民共和国教育部和国家语言文字工作委员会开始就修改后的简体字表征询公众意见。[4][5][6][7]《通用规范汉字表》于2013年9月正式实施，[8]收录2574个简化字[8]。
历史
主条目：汉字简化 § 中国大陆

1920年，钱玄同在《新青年》上发表了《减省汉字笔画的提议》。1931年，“中国新文字第一次代表大会”在海参崴举行，中国共产党代表瞿秋白、吴玉章等人与苏联共同草拟“北方话拉丁化新文字”，并发表了中苏双方13条共同宣言，其中强调“中国汉字是古代封建社会的产物，成了统治阶级压迫劳苦群众的工具之一”， 因此“要根本废除象形文字，以纯粹的拼音文字代替”，并“反对用象形文字的笔划来拼音或注音”。宣言还明确“反对中国资产阶级的所谓统一国语运动”。 [9]

1923年，胡适在《国语月刊·汉字改革号》的《卷头言》中说：“这二千年的中国的小百姓不但做了很惊人的文法革新，他们还做了一件同样惊人的革新事业：就是汉字形体上的大改革，就是‘破体字’的创造与提倡。”[10]

1934年12月，鲁迅在《关于新文字》一文中写道：“方块字真是愚民政策的利器……汉字也是中国劳苦大众身上的一个结核，病菌都潜伏在里面，倘不首先除去它，结果只有自己死。”鲁迅于1936年10月更坚决地说：“汉字不灭，中国必亡。”理由是：“因为汉字的艰深，使全中国大多数的人民，永远和前进的文化隔离，中国的人民，决不会聪明起来，理解自身所遭受的压榨，理解整个民族的危机。”[11]

到了1935年，中华民国教育部发布了第11400号部令，公布了《第一批简体字表》。然而，仅仅一年之后，该方案就因为遭到中国国民党元老戴季陶的强烈反对而作罢。[12]戴季陶称推行简体字是比亡国灭种还可怕的事。[13]中华民国教育部在行政院的命令下，训令暂不推行简体字。[12][14]与此同时在解放区，群众自发采用并创造了许多简体字，例如“辽、远、拥、护、开、关、运、动、奋、斗、敌、卫、胜、艺、习、团”等，这些字习惯上被叫作“解放字”。 [15][16]

中华人民共和国成立后，简化汉字的进程加速推进。1952年，在中央教育部的筹划下，中国文字改革研究委员会成立。中国科学院语言研究所俄籍顾问谢尔应琴柯（Selchiuchinko）在该委员会会议上指示：“严格的拼音原则是采用中国共产党员在苏联创制的那套拼音文字是最合理的”。[17]1953年，文字改革研究委员会汉字整理组成立，开始着手拟定《常用汉字简化表草案》，专家们以普遍通行的简体字为主，辅以草书楷化的方法，选定了700个简体字，拟出第一稿。 但是，毛泽东看过后并不满意，他指出：“700个简体字还不够简，作简体字要多利用草体，找出简化规律，作出基本形体，有规律地进行简化。汉字的数量也必须大大减缩，一个字可以代替好几个字，只有从形体上和数量上同时精简，才算得上简化。”[18]由于汉字难以在短时间内完全改为拼音文字，文字改革委员会预计会有一段新旧文字并用的过渡时期，因此决定通过破坏汉字形体、削弱汉字效用，以简化字创造汉字内部利于拼音化的条件。 [19]

1954年12月，中国文字改革研究委员会改组为中国文字改革委员会（简称文改会），直属于国务院。随后，文改会在1955年公布了《汉字简化方案草案》，其中包含三个表：〈798个汉字简化表〉、〈拟废除的400个异体字表〉、〈汉字偏旁手写简化表〉。其中第二表就是其后分拆出的《第一批异体字整理表》的草案。 〈汉字简化方案草案说明〉中指出：“通过这个草案的讨论，我们希望大家共同努力研究中国文字改革问题，为今后进一步整理汉字和实行拼音文字创造更有利的条件。” 同年，文改会在《汉字简化方案草案说明》中再次强调汉字难以在短时间内完全改为拼音文字，即使在开始实行拼音文字后，仍会有一个新旧文字并用的过渡时期，汉字仍然是一定时期内必须使用的重要工具，并提出了三种简化汉字的方法：笔画、字数、写法。 [20]1956年1月28日，中华人民共和国国务院全体会议第23次会议通过《关于公布汉字简化方案的决议》，并将《汉字简化方案》全文发表在1月31日的《人民日报》上。 [18]然而，汉字简化方案的推行并非一帆风顺，一些人对这项改革提出了质疑。1958年，周恩来发表了《当前文字改革的任务》报告，指责“一些右派份子对文字改革进行了恶毒的攻击，说汉字简化搞糟了，群众都反对，要国务院收回成命，把《汉字简化方案》撤回”。 他表示汉字简化“是符合群众利益并且受到群众热烈欢迎的好事”，“应该给以坚决支持”。

1964年5月，中国文字改革委员会出版《简化字总表》，该表共收录2236字，其中第一表是352个不作偏旁使用的简化字，第二表是132个可作简化偏旁的简化字，第三表是由第二表类推的1754字。[18]自1976年起，新加坡教育部发布的汉字简体字表也与中国大陆的简化字完全相同。[21] 

\end{document}